% Options for packages loaded elsewhere
\PassOptionsToPackage{unicode}{hyperref}
\PassOptionsToPackage{hyphens}{url}
\documentclass[
  ignorenonframetext,
]{beamer}
\newif\ifbibliography
\usepackage{pgfpages}
\setbeamertemplate{caption}[numbered]
\setbeamertemplate{caption label separator}{: }
\setbeamercolor{caption name}{fg=normal text.fg}
\beamertemplatenavigationsymbolsempty
% remove section numbering
\setbeamertemplate{part page}{
  \centering
  \begin{beamercolorbox}[sep=16pt,center]{part title}
    \usebeamerfont{part title}\insertpart\par
  \end{beamercolorbox}
}
\setbeamertemplate{section page}{
  \centering
  \begin{beamercolorbox}[sep=12pt,center]{section title}
    \usebeamerfont{section title}\insertsection\par
  \end{beamercolorbox}
}
\setbeamertemplate{subsection page}{
  \centering
  \begin{beamercolorbox}[sep=8pt,center]{subsection title}
    \usebeamerfont{subsection title}\insertsubsection\par
  \end{beamercolorbox}
}
% Prevent slide breaks in the middle of a paragraph
\widowpenalties 1 10000
\raggedbottom
\AtBeginPart{
  \frame{\partpage}
}
\AtBeginSection{
  \ifbibliography
  \else
    \frame{\sectionpage}
  \fi
}
\AtBeginSubsection{
  \frame{\subsectionpage}
}
\usepackage{iftex}
\ifPDFTeX
  \usepackage[T1]{fontenc}
  \usepackage[utf8]{inputenc}
  \usepackage{textcomp} % provide euro and other symbols
\else % if luatex or xetex
  \usepackage{unicode-math} % this also loads fontspec
  \defaultfontfeatures{Scale=MatchLowercase}
  \defaultfontfeatures[\rmfamily]{Ligatures=TeX,Scale=1}
\fi
\usepackage{lmodern}
\ifPDFTeX\else
  % xetex/luatex font selection
\fi
% Use upquote if available, for straight quotes in verbatim environments
\IfFileExists{upquote.sty}{\usepackage{upquote}}{}
\IfFileExists{microtype.sty}{% use microtype if available
  \usepackage[]{microtype}
  \UseMicrotypeSet[protrusion]{basicmath} % disable protrusion for tt fonts
}{}
\makeatletter
\@ifundefined{KOMAClassName}{% if non-KOMA class
  \IfFileExists{parskip.sty}{%
    \usepackage{parskip}
  }{% else
    \setlength{\parindent}{0pt}
    \setlength{\parskip}{6pt plus 2pt minus 1pt}}
}{% if KOMA class
  \KOMAoptions{parskip=half}}
\makeatother
\usepackage{longtable,booktabs,array}
\usepackage{caption}
\captionsetup[table]{skip=6pt}
\newcounter{none} % for unnumbered tables
\usepackage{calc} % for calculating minipage widths
\usepackage{caption}
% Make caption package work with longtable
\makeatletter
\def\fnum@table{\tablename~\thetable}
\makeatother
\setlength{\emergencystretch}{3em} % prevent overfull lines
\providecommand{\tightlist}{%
  \setlength{\itemsep}{0pt}\setlength{\parskip}{0pt}}
% Manuscript macro/package fallbacks (newcommand -> providecommand).
% Core mathematical packages
\usepackage{amsmath,amssymb,amsthm}
\usepackage{mathtools}
% Keep long deployment prose and references within the text block.
% Algorithm formatting
% Tables
\usepackage{booktabs}
\usepackage{multirow}
% Graphics and floats
% Typography and layout
% Page layout (slightly smaller margins for academic journal style)
% Running headers/footers
% Caption formatting
% Colored boxes for callouts
% Cross-referencing with smart naming
% Hyperlink styling
% Configure cleveref for custom environments
% Theorem environments with consistent numbering
% Code listings and raw implementation examples in the manuscript
\usepackage{listings}
\lstset{basicstyle=\ttfamily\small,breaklines=true,columns=fullflexible}
% Math operators
\DeclareMathOperator*{\argmax}{arg\,max}
\DeclareMathOperator*{\argmin}{arg\,min}
\DeclareMathOperator{\sign}{sign}
\DeclareMathOperator{\KL}{KL}
\DeclareMathOperator{\Tr}{Tr}
% Custom commands for notation consistency
\providecommand{\calA}{\mathcal{A}}
\providecommand{\calB}{\mathcal{B}}
\providecommand{\calC}{\mathcal{C}}
\providecommand{\calD}{\mathcal{D}}
\providecommand{\calF}{\mathcal{F}}
\providecommand{\calG}{\mathcal{G}}
\providecommand{\calH}{\mathcal{H}}
\providecommand{\calI}{\mathcal{I}}
\providecommand{\calM}{\mathcal{M}}
\providecommand{\calO}{\mathcal{O}}
\providecommand{\calP}{\mathcal{P}}
\providecommand{\calS}{\mathcal{S}}
\providecommand{\calT}{\mathcal{T}}
\providecommand{\calW}{\mathcal{W}}
% Note: \Phi is a standard LaTeX command, do not redefine
% Trust notation
\providecommand{\trust}[2]{\mathcal{T}_{#1 \to #2}}
\providecommand{\trustt}[3]{\mathcal{T}_{#1 \to #2}^{#3}}
% Belief notation
\providecommand{\belief}[2]{\mathcal{B}_{#1}(#2)}
\providecommand{\belieft}[3]{\mathcal{B}_{#1}^{#2}(#3)}
% Cognitive state
\providecommand{\cogstate}[1]{\sigma_{#1}}
% Attack notation
\providecommand{\attack}[1]{\mathcal{A}_{#1}}
\providecommand{\adversary}[1]{\Omega_{#1}}
% Defense notation
\providecommand{\firewall}{\mathcal{F}}
\providecommand{\sandbox}{\mathcal{S}_{box}}
% Probability and expectation
\providecommand{\E}{\mathbb{E}}
\providecommand{\Prob}{\mathbb{P}}
\providecommand{\indicator}{\mathbb{1}}
% QED symbol
\providecommand{\qedsymbol}{$\blacksquare$}
% Front matter opens on the cover/title page (no list of figures/tables precedes it).

% Slide layout defaults for warning-clean scientific decks.
\usepackage{etoolbox}
\IfFileExists{xurl.sty}{\usepackage{xurl}}{}
\IfFileExists{seqsplit.sty}{\usepackage{seqsplit}}{\newcommand{\seqsplit}[1]{#1}}
\protected\def\breakseq#1{\seqsplit{#1}}
\protected\def\breaktt#1{\begingroup\ttfamily\seqsplit{#1}\endgroup}
\setlength{\emergencystretch}{6em}
\tolerance=5000
\hbadness=10000
\hfuzz=1pt
\setlength{\tabcolsep}{2pt}
\AtBeginEnvironment{longtable}{\tiny\renewcommand{\arraystretch}{0.86}\setlength{\tabcolsep}{1pt}}
\AtBeginEnvironment{tabular}{\tiny\renewcommand{\arraystretch}{0.86}\setlength{\tabcolsep}{1pt}}
\AtBeginEnvironment{equation}{\tiny}
\AtBeginEnvironment{equation*}{\tiny}
\AtBeginEnvironment{align}{\tiny}
\AtBeginEnvironment{align*}{\tiny}
\AtBeginEnvironment{itemize}{\footnotesize}
\AtBeginEnvironment{enumerate}{\footnotesize}
\AtBeginEnvironment{description}{\footnotesize}
\setbeamerfont{caption}{size=\tiny}
\setbeamerfont{caption name}{size=\tiny}
\setbeamerfont{normal text}{size=\small}
\setbeamerfont{frametitle}{size=\small}
\setbeamerfont{section title}{size=\footnotesize}
\setbeamerfont{subsection title}{size=\footnotesize}
\setbeamertemplate{section page}{%
  \centering
  \begin{beamercolorbox}[sep=12pt,center,wd=\paperwidth]{section title}
    \parbox{0.86\paperwidth}{\centering\usebeamerfont{section title}\insertsection\par}
  \end{beamercolorbox}
}
\setbeamertemplate{subsection page}{%
  \centering
  \begin{beamercolorbox}[sep=8pt,center,wd=\paperwidth]{subsection title}
    \parbox{0.86\paperwidth}{\centering\usebeamerfont{subsection title}\insertsubsection\par}
  \end{beamercolorbox}
}
\setlength{\abovecaptionskip}{2pt}
\setlength{\belowcaptionskip}{0pt}

% Natbib and cross-reference fallbacks — slides don't load natbib
% or cleveref, but combined-PDF manuscript prose may emit these
% commands. The fallback renders citations as a bracketed key list
% and cross-references as detokenized labels so slides don't fail on
% undefined control sequences or raw underscores. \providecommand is
% a no-op if packages load later.
\providecommand{\citep}[1]{[#1]}
\providecommand{\citet}[1]{#1}
\providecommand{\citealp}[1]{#1}
\providecommand{\citeauthor}[1]{#1}
\providecommand{\citeyear}[1]{#1}
\providecommand{\cref}[1]{\texttt{\detokenize{#1}}}
\providecommand{\Cref}[1]{\texttt{\detokenize{#1}}}
\providecommand{\crefrange}[2]{\texttt{\detokenize{#1}}--\texttt{\detokenize{#2}}}
\providecommand{\Crefrange}[2]{\texttt{\detokenize{#1}}--\texttt{\detokenize{#2}}}
\providecommand{\paragraph}[1]{\textbf{#1}\ }

% Environment-providing packages carried over from the manuscript.
\usepackage{algpseudocode}

% Non-floating stand-in for the `algorithm` float.
\newenvironment{algorithm}[1][]{%
  \par\medskip\noindent\rule{\linewidth}{0.4pt}\par\nobreak\small
  \renewcommand{\caption}[1]{\par\noindent\textbf{##1}\par}}{%
  \par\nobreak\noindent\rule{\linewidth}{0.4pt}\par\medskip}

\newtheorem{proposition}{Proposition}
\newtheorem{hypothesis}{Hypothesis}
\newtheorem{remark}[theorem]{Remark}
\newtheorem{axiom}{Axiom}
\newtheorem{property}{Property}
\usepackage{bookmark}
\IfFileExists{xurl.sty}{\usepackage{xurl}}{} % add URL line breaks if available
\urlstyle{same}
% fallback for those not using the hyperref driver hyperxmp:
\makeatletter
\@ifundefined{xmpquote}{\newcommand{\xmpquote}[1]{#1}}{}
\makeatother
\hypersetup{
  hidelinks,
  pdfcreator={LaTeX via pandoc}}

\author{\texorpdfstring{}{}}
\date{}

\begin{document}

\begin{frame}
\newpage
\end{frame}

\section{Where We Stand: A Call to Build}\label{sec:conclusion}

\begin{frame}[allowframebreaks]{Where We Stand: A Call to Build}
This series began with a theory (Part 1) and moved to an experiment
(Part 2). It ends here, with a synthesis and a call to engineering.
\end{frame}

\subsection{The Theory Holds}\label{the-theory-holds}

\begin{frame}[allowframebreaks]{The Theory Holds}
We proved that trust can be bounded. We proved that defenses can be
composed algebraically. We proved that stealth and impact are inversely
related. These are not just academic curiosities; they are foundational
constraints for secure cognitive systems.

The theoretical foundations have deepened since the first drafts of this
guide. Three new formal results from Part 2 strengthen the case:

\textbf{Categorical guarantee}: Theorems CT.1--CT.3 show defense
composition in CIF is constrained by the DefenseCategory structure: a
detection-preserving chain cannot yield a non-detecting composite under
the stated assumptions. That is a property of the composition law, not a
separate empirical fit.

\end{frame}

\begin{frame}[allowframebreaks]{The Theory Holds}
\textbf{Free energy connection}: FEP.1--FEP.2 state a correspondence
between CIF's trust update and precision-weighted active inference under
the generative model used in Part\textasciitilde2, so variational free
energy gives one interpretive lens for why decay and sandboxing change
belief updates the way they do.

\textbf{Geometric bound}: Theorem CG.1 establishes that the belief
sandbox imposes a hard geodesic boundary on belief manipulation: no
attack can move an agent's beliefs beyond the Riemannian radius \(\rho\)
without triggering the sandbox. This is a structural guarantee
independent of attack sophistication---it holds for any manipulation
that preserves probability mass, including attacks that current
classifiers cannot detect by pattern.
\end{frame}

\begin{frame}[fragile,allowframebreaks]{The Code Works---At Three
Levels}
\protect\phantomsection\label{the-code-worksat-three-levels}
Saying ``the code works'' requires specifying what that means. There are
three honest answers, each true at a different level:

\textbf{Level 1 --- The architecture is correctly implemented}
(confidence: high, conditional on the current project test run). Every
defense module---firewall, sandbox, trust calculus, tripwires, Byzantine
consensus, provenance, drift detection, invariant checker---is
independently tested. The \texttt{SeriesPipeline} routes the 950-attack
corpus through all 8 modules with 0\% routing failure in the reported
Part 2 run. Use the current \texttt{uv\ run\ pytest} output as the
source of truth for test counts and pass rate.

\end{frame}

\begin{frame}[fragile,allowframebreaks]{The Code Works---At Three
Levels}
\textbf{Level 2 --- The pipeline detects attacks in practice}
(confidence: moderate). The multi-seed pipeline analysis (30 seeds,
Claude Code architecture) achieves a mean detection rate of 44.8\%
{[}95\% CI: 43.2\%, 46.4\%{]}. The LLM-backed multiagent validation
(\(N=10\), Gemma 3 4B) achieves 80\%--100\% across two architectures.
These are meaningful but not high detection rates---they reflect adapter
implementation at CMMI Level 3 (Statistical), not the design ceiling.

\end{frame}

\begin{frame}[fragile,allowframebreaks]{The Code Works---At Three
Levels}
\textbf{Level 3 --- The defense ceiling is achievable} (confidence:
moderate-high). The parametric simulation (\(N=3{,}800\)) establishes
that fully-mature (Level 5) adapters achieve 96--100\% detection,
consistent with the formal design. The gap between Level 3 (44.8\%) and
Level 5 (96\%) is an engineering challenge, not a theoretical
limitation. The roadmap in Part 2 projects +35--41 percentage points of
improvement through adapter maturation.

\end{frame}

\begin{frame}[fragile,allowframebreaks]{The Code Works---At Three
Levels}
The honest operational posture is Level 2: deploy CIF for meaningful
protection against \(\Omega_1\)--\(\Omega_3\) attacks today, while
investing in adapter maturation for \(\Omega_4\)--\(\Omega_5\) coverage.
Do not rely on 94\% detection for life-safety applications until your
adapters reach Level 4--5 and have been validated against your threat
model.
\end{frame}

\subsection{Summary of Practical
Recommendations}\label{summary-of-practical-recommendations}

\begin{frame}[allowframebreaks]{Summary of Practical Recommendations}
The preceding sections distill the CIF series into actionable guidance.
The core recommendations are:

\end{frame}

\begin{frame}[allowframebreaks]{Summary of Practical Recommendations}
\begin{enumerate}
\item
  \textbf{Adopt layered defense from the start} (Pitfall 2), then
  measure what each layer is worth. In the parametric model the full CIF
  stack reached a 96--100\% parametric detection ceiling that no partial
  configuration matched. The real pipeline teaches the sharper lesson:
  on Part\textasciitilde2's 98-attack ablation corpus the Invariants
  checker alone reaches 92.9\% against the full stack's 95.9\%, and five
  of the eight modules show no measurable marginal contribution. Design
  security into the architecture rather than bolting it on, but do not
  assume every layer you add is earning its latency --- check each one's
  marginal contribution against your own traffic.
\item
  \textbf{Implement trust decay on every delegation chain} (Pitfall 1).
  The Trust Calculus with \(\delta \leq 0.8\) prevents trust laundering
  across all tested architectures. This is not optional hardening---it
  is the structural foundation that prevents systemic compromise from
  local failures.
\item
  \textbf{Deploy tripwires with rotation} (Pitfall 5). Identity,
  boundary, and principal canaries provide continuous detection of
  belief manipulation. Static tripwires degrade over time; automated
  rotation maintains detection coverage.
\item
  \textbf{Monitor for progressive drift, not just sudden changes}
  (Pitfall 6). Sliding-window drift detection catches the gradual belief
  corruption that per-update thresholds miss. Track cumulative
  deviation, not just per-step delta.
\item
  \textbf{Log everything} (Pitfall 7). Full belief provenance,
  inter-agent message history, and cognitive state snapshots are
  essential for post-incident forensics. Without them, attack
  reconstruction is impossible.
\end{enumerate}
\end{frame}

\subsection{Maturity Roadmap}\label{maturity-roadmap}

\begin{frame}[allowframebreaks]{Maturity Roadmap}
Organizations adopting cognitive security should plan a staged
deployment aligned with the three profiles evaluated in Section 5:

\textbf{Stage 1: Minimal Viable Implementation} (Weeks 1--4)

\end{frame}

\begin{frame}[allowframebreaks]{Maturity Roadmap}
\begin{itemize}
\tightlist
\item
  Deploy the Cognitive Firewall at all agent ingress points
\item
  Set trust decay \(\delta = 0.80\) on all delegation chains
\item
  Place at least one identity tripwire per agent
\item
  Expected outcome: Low-effort attacks reduced from 100\% baseline
  success to \textless5\%
\end{itemize}

\end{frame}

\begin{frame}[allowframebreaks]{Maturity Roadmap}
\textbf{Stage 2: Balanced Deployment} (Months 2--3)

\end{frame}

\begin{frame}[allowframebreaks]{Maturity Roadmap}
\begin{itemize}
\tightlist
\item
  Add Belief Sandboxing with \(\kappa = 2\) corroboration
\item
  Deploy drift detection with sliding-window analysis
\item
  Implement structured provenance logging
\item
  Tune thresholds against representative attack samples from Part 2's
  corpus
\item
  Expected outcome: a 96--100\% parametric detection ceiling at
  \textasciitilde20\% latency overhead; real-pipeline performance must
  be measured separately
\end{itemize}

\end{frame}

\begin{frame}[allowframebreaks]{Maturity Roadmap}
\textbf{Stage 3: High Assurance} (Months 4--6)

\end{frame}

\begin{frame}[allowframebreaks]{Maturity Roadmap}
\begin{itemize}
\tightlist
\item
  Enable Byzantine consensus for critical collective decisions
  (\(n \geq 3f + 1\))
\item
  Implement aggressive trust decay (\(\delta = 0.60\)) for autonomous
  operations
\item
  Deploy orchestrator-specific monitoring (Pitfall 8)
\item
  Establish canary rotation schedules and drift baseline refresh cycles
\item
  Expected outcome: 95--98\% detection across all categories, suitable
  for unsupervised autonomous agents
\end{itemize}

\end{frame}

\begin{frame}[allowframebreaks]{Maturity Roadmap}
\textbf{Stage 4 (Months 7--12): Adapter Maturation and Verified
Coverage} \emph{Target: 80--86\% detection rate, Level-4--5 adapters for
primary attack categories}

\end{frame}

\begin{frame}[allowframebreaks]{Maturity Roadmap}
\begin{itemize}
\tightlist
\item
  Advance Cognitive Firewall adapter from Level 3 (Statistical) to Level
  4 (Adaptive) via online learning against observed attack distributions
  (+15--20 pp DR)
\item
  Advance Trust Calculus adapter from Level 3 to Level 4 via dynamic δ
  calibration from operational trust logs (+5--8 pp DR)
\item
  Deploy collective free energy monitoring for emergent misalignment
  (Direction 5): target 70\% detection on colony-scale drift scenarios
\item
  Conduct red team exercise with \(N \geq 50\) attacks per category for
  statistically powered evaluation (see Part\textasciitilde2's
  power-analysis table)
\item
  Validate parametric-to-empirical closure: confirm empirical DR
  approaches parametric ceiling as adapters mature
\end{itemize}

\end{frame}

\begin{frame}[allowframebreaks]{Maturity Roadmap}
\emph{Milestone}: 80--86\% empirical detection rate on the standard
attack corpus with 95\% HDI width \(\leq 10\%\).

At each stage, validate against the Summary Checklist in Section 6
before advancing. Address unchecked items before production deployment.
\end{frame}

\subsection{The Next Step is Yours}\label{the-next-step-is-yours}

\begin{frame}[allowframebreaks]{The Next Step is Yours}
As we move beyond simple prompt engineering, the era of
\textbf{Cognitive Security Engineering} has begun.

We are no longer just asking LLMs to write poems; we are asking them to
run the world. If we want them to do that safely, we cannot rely on luck
or better fine-tuning. We need structure. We need rigorous,
mathematically grounded, architecturally sound defense systems.

The CIF is our contribution to that structure. It is a toolbox, not a
bible. Take it. Fork it. Break it. Improve it.

\end{frame}

\begin{frame}[allowframebreaks]{The Next Step is Yours}
The agents are coming online. Let's make sure they are safe.
\end{frame}

\subsection{A Note on Uncertainty}\label{sec:uncertainty-note}

\begin{frame}[allowframebreaks]{A Note on Uncertainty}
This guide has tried to be honest about what CIF can and cannot do. The
practical guidance---configuration tables, playbooks, monitoring
thresholds, case studies---is grounded in the best available evidence.
But that evidence has real limitations that operators should carry
forward into their deployment decisions:

\textbf{Sample sizes are small}. The LLM validation used \(N=5\) attacks
per architecture. The colony benchmarks used 1 scenario per attack type.
The multi-seed analysis used 30 seeds. These are sufficient for
preliminary evidence but severely underpowered for precise estimation.
Required sample sizes for \(\pm 5\%\) precision are \(N \geq 246\) per
evaluation mode. Treat all reported detection rates as estimates with
wide uncertainty, not precise measurements.

\end{frame}

\begin{frame}[allowframebreaks]{A Note on Uncertainty}
\textbf{The gap is real at one end}. The 4--51 percentage-point gap
between the parametric ceiling (96--100\%) and the real pipeline spans
two unlike measurements. The wide end is the one the Bayes factor speaks
to: the distance from the ceiling to the 30-seed pipeline mean of
44.8\%, which is not noise---the Bayes factor for a true performance gap
exceeds \(10^6\) (decisive evidence). The narrow end is the distance
from the ceiling to the 95.9\% the pipeline reaches on the 98-attack
ablation corpus after the Invariants rewrite; no Bayes factor has been
computed for that arm, and it should not borrow the first one's
authority. The wide end is what reflects adapter implementation
maturity, and maturation takes time and resources. Plan your deployment
timeline and security posture accordingly.

\end{frame}

\begin{frame}[allowframebreaks]{A Note on Uncertainty}
\textbf{Adaptive adversaries are not modeled}. All evaluations used a
fixed attack corpus. A sophisticated adversary who observes CIF's
defenses and adapts (Debenedetti et al.'s adaptive attacks
\cite{adaptive2025attacks}) could achieve lower detection rates than
reported. The game-theoretic analysis (Part 2) establishes the Nash
equilibrium for the current payoff matrix, but a patient adversary will
probe for gaps. The layered defense architecture provides
resilience---bypassing one layer still encounters others---but no
detection system is perfectly robust to adaptive adversaries.

\end{frame}

\begin{frame}[allowframebreaks]{A Note on Uncertainty}
\textbf{Uncertainty is not a reason not to deploy}. CIF provides
meaningful, formally-grounded protection that no single-agent guardrail
system provides. The trust calculus's structural guarantee (trust cannot
be amplified through delegation, regardless of detection rates) is
unconditional. The compositional algebra ensures that layered defenses
provide more coverage than any single mechanism. Deploy CIF with honest
expectations, monitor carefully, and improve iteratively. That is sound
engineering practice---not a confession of inadequacy.
\end{frame}

\end{document}
